\centerline{\bf \Huge Сравнение по модулю. Повторение}
\title{ Сравнение по модулю. Повторение}

\problem{1} Найдите все возможные остатки от деления квадрата целого числа:

а) на 3; б) на 4; в) на 5.

\noindent
Объясните, почему число 100 $\dotsc$ 04 не может быть квадратом целого числа

\noindent\textit{Определение.} Запись $a \equiv b \ (mod \ m)$ означает, что числа $a$ и $b$ имеют одинаковые остатки при делении на $m$. Читается так: $a$ сравнимо с $b$ по модулю $m$.

\problem{2} Докажите, что $a \equiv b \ (mod \ m)$ тогда и только тогда, когда $a - b$ делится на $m$.

\problem{3} \textit{(Свойства сравнений)} Пусть $a \equiv b \ (mod \ m)$, $c \equiv d \ (mod \ m)$. Докажите, что:

• $a + c \equiv b + d \ (mod \ m)$ — сравнения по одному и тому же модулю можно складывать друг с другом;

• $ac \equiv bd \ (mod \ m) -$ сравнения по одному и тому же модулю можно умножать друг на друга;

• $a^n \equiv b^n \ (mod \ m)$, $n \in \mathbb{N}$ — сравнение можно возвести в любую натуральную степень;

• $ka \equiv kb \ (mod \ m), k \in \mathbb{Z} $— сравнение можно умножить на любое целое число.

\problem{4} Найдите остаток от деления: а) $5^{20}$ на $24;$ б) $3^{66}$ на $28.$

\problem{5} Докажите, что $16^{2014} + 33^{2015}$ делится на $17$.

\problem{6} При каких натуральных $n$ число $2^n - 1$ делится на $7$?

\problem{7} Докажите, что при любом натуральном $n$ число ${37}^{n+2} + {16}^{n+1} + {23}^n$ делится на $7$.

\problem{8} Докажите, что для $n$-значного числа $\overline{a_1a_2\ldots a_n}$ верно:\\
а) $a_1a_2 . . . a_n \equiv a_1 + a_2 + . . . + a_n \ (mod \ 9);$\\
б)$a_1a_2 . . . a_{n-1}a_n \equiv a_n - a_{n-1} + a_{n-2} - . . . + (-1)^{n-1}a_1 \ (mod \ 11).$\\
Выведите отсюда признаки деления на 9 и 11.

\newpage

\textbf{Задачи для самопроверки:}

\problem{1} Делится ли число $21^{10} - 1$ на $2200?$

\problem{2} Найти последнюю цифру числа  $1\cdot2 + 2\cdot3 + ... + 999\cdot1000.$

\problem{3} Найти остаток  $(116 + 17^{17})^{21}·7^{49}$  от деления на 8.

\problem{4} Даны 111 различных натуральных чисел, не превосходящих 500.
Могло ли оказаться, что для каждого из этих чисел его последняя цифра совпадает с последней цифрой суммы всех остальных чисел?
